\section{Source audit, proposed corrections, and integration decisions}
\label{sec:audit}

The source comparison did not uncover a false theorem in the particular
latest Euler, Lerch, and full-Cayley sections used as premises here.
It did reveal several places where reading an older fragment in isolation
would give an incorrect impression of the project's current status.
The following proposed actions distinguish a newly proved conjecture,
an extension of a correct theorem, and an editorial status update.

\subsection{Concrete source actions}

\begin{longtable}{@{}>{\raggedright\arraybackslash}p{0.23\textwidth}>{\raggedright\arraybackslash}p{0.32\textwidth}>{\raggedright\arraybackslash}p{0.39\textwidth}@{}}
\toprule
Source item & Status at the pinned baseline & Proposed action\\
\midrule
\endfirsthead
\toprule
Source item & Status at the pinned baseline & Proposed action\\
\midrule
\endhead
Near-critical Euler crossing & Explicit conjecture
\texttt{research:conj:crossing} in \cite{RepoEuler}. & Replace the conjecture
by Theorem~\ref{eutrans:thm:crossing}; retain the compact positive
$b$-range and fixed derivative order in its statement.\\[5pt]
Lerch sign-block question & The latest report gives a sufficient bound
and asks for the least length \cite{RepoLerch}. & Insert
Theorem~\ref{newdiag:thm:sharp-block}. The bound is optimal; the reciprocal
angle case has length $\pi/\theta+1$.\\[5pt]
Higher critical sums & The first weighted sum is proved; higher sums
are proposed in \cite{RepoLerch}. & Insert
Theorem~\ref{newdiag:thm:finiteparts} and the explicit $p=2$ formula.
Retain the conjugate-singularity subtractions.\\[5pt]
Full Cayley ideal & A full projector, all-depth dimensions, and the
frozen target's nonmembership are already proved
\cite{depth:PriorProjection}. & Add the depth-preserving retraction
and filtered counts. Keep $\Pi C=\Pi$ distinct from $C\Pi=\Pi$.\\[5pt]
Classical Stieltjes indices $6,7$ & Complete all-derivative transition
counts are already in the manuscript \cite{RepoZeros}. & Mark any
older first-derivative conjectures as superseded. Do not present those
counts as a new result of this report.\\[5pt]
$S_4$, $S_6$, and current $S_8$ & $S_4$ is proved in later material;
the current $S_6,S_8$ candidates remain conjectural
\cite{RepoManuscript}. & Preserve these statuses. The 30-term normal form
is a formal reduction of the $S_6$ target, not a period certificate.\\
\bottomrule
\end{longtable}

For clarity, the already established Stieltjes result is
\[
 \#\{a>0:\gamma_6^{(k)}(a)=0\}=
 \begin{cases}4,&1\leq k\leq3,\\6,&k\geq4,\end{cases}
 \qquad
 \#\{a>0:\gamma_7^{(k)}(a)=0\}=
 \begin{cases}5,&1\leq k\leq4,\\7,&k\geq5.
 \end{cases}
\]
All these zeros are simple, with the global count justified by the
source's analytic bound and rational sign certificates. This report does
not replay those inherited certificates and does not claim their
discovery. Likewise, the dominant inverse pair for the Lerch velocities
is an inherited selected-sheet theorem, not an open assumption.

\subsection{Proof obligations that should remain explicit}

\paragraph{Joint asymptotics require a joint remainder.}
The fixed-parameter Euler density expansion is correct within its stated
scope. Substituting $a=1-\varepsilon$ into a remainder whose constant can
depend on $a$ would not prove the conjectured crossing law. The algebraic
coefficient tends to zero and the seed survives. The new
Lemma~\ref{eutrans:lem:density} supplies the missing $\varepsilon$ factor;
Lemma~\ref{eutrans:lem:joint} controls both competing scales separately.
The proposed correction is to any attempted use beyond the old theorem's
scope, rather than to that theorem itself.

\paragraph{Sharp finite parts see both dominant singularities.}
The first weighted Lerch sum has no divergent conjugate contribution.
At the next order the conjugate term has magnitude $\sqrt N$.
Consequently a proposed higher-order finite part subtracting only the
nonoscillatory powers cannot converge in general. Equation
\eqref{newdiag:eq:E} is the required addition. This is a concrete
correction to the naive extension suggested by the first-order pattern;
it is not a claim that the source stated that extension as a theorem.

\paragraph{Depth is attached to a specified relation system.}
The quotient in Section~\ref{depthcay:sec} imposes the full Cayley ideal,
and its depth minima are proved within that quotient. Extra shuffle,
stuffle, distribution, functional, or period relations can identify
additional elements. Neither a formal dimension nor an unsuccessful
bounded search gives a lower bound for the dimension of numerical
periods. The source ledger already emphasizes this distinction
\cite{RepoManuscript}; the proposed integration preserves it.

\paragraph{Translation and endpoint conventions matter.}
In the Gamma energy the shift is periodic, with a break at $1-h$.
Replacing it by $\log\Gamma(x+h)-\log\Gamma(x)$ over the whole interval
changes the problem. For complex spectral parameters the symmetric
polylogarithm pair in \eqref{eq:polylog-master} must not be replaced
by a real part. The removable pole in
$u(u-1)\zeta(1+u)$ must be canceled before order differentiation.
Finally, repeated definite primitives require the polynomial subtraction
in \eqref{eq:shift-definite}; an indefinite primitive alone does not fix
the endpoint data.

\subsection{Suggested placement}

The translation and Gamma sections fit the manuscript's integrated
zeta-jet and Stieltjes-antiderivative chapter, with the sharp regularity
result cross-referenced from the spectral discussion. The uniform Euler
section extends the real-order kernel and fractional continuation.
The sign-block and critical-sum results belong immediately after the
dominant Lambert-pair theorem. The Cayley section belongs after the full
ideal projector, followed by the weight-seven search discussion.
Complete source paths and theorem labels, without dependence on changing
chapter numbers, are in \texttt{integration/claim\_ledger.json} and
\texttt{integration/INTEGRATION.md}.

\section{Further research questions and proposed work}
\label{sec:questions}

The following questions are deliberately stated with a target and a
specific obstacle. None is used in the proofs above. The first group
extends the new analytic identities; the later groups connect them back
to the remaining arithmetic and geometric questions in the source.

\subsection{Translation identities and the Stieltjes hierarchy}

\begin{question}[Rational shifts and explicit character coordinates]
For $h=p/q$, produce a systematic normal form for the translated
pairings in terms of Dirichlet $L$-function order derivatives and
conductor logarithms. Which low-conductor specializations simplify
to small explicit Gamma, polygamma, and zeta identities?
\end{question}
The exact finite Fourier identity
\[
 \Li_s(e^{2\pi ip/q})=q^{-s}\sum_{a=1}^q
 e^{2\pi ipa/q}\zeta(s,a/q),\qquad \Re s>1,
\]
and its order derivatives already provide a finite conversion.
The next step is to combine it with the repository's conductor and
distribution reductions, retaining the logarithmic terms generated
by $q^{-s}$. This is a coordinate problem; it should not be confused
with proving that the resulting constants are independent.

\begin{question}[Higher translation energies]
For $r\geq2$, determine the extrema and concavity intervals of
$V_r(h)=\int_0^1[\Delta_h\zeta^{(r)}(0,x)]^2\,dx$.
Does the half-period remain the unique global maximizer in every order?
\end{question}
Symmetry ensures a stationary point at $1/2$, but a positive cosine-series
coefficient does not by itself force that point to be the global maximum.
The master identity gives the complete derivative formula. A proof
would need a suitable sign theorem for combinations of higher
Stieltjes constants, extending the elementary $r=1$ argument.

\begin{question}[Exact endpoint Gram determinants]
Find closed formulas, recurrence relations, or asymptotic bounds for
the determinants of the endpoint-polynomial Gram matrices, and for
the orthogonal polynomials associated with their two-component
integral representation.
\end{question}
Their dependence on the logarithmic shift is triangular and their
positivity is exact. The unexplained part is the all-order determinant
and orthogonalization structure, whose entries contain zeta values.
An explicit recurrence could also improve high-order evaluation of
the Stieltjes translation covariances.

\begin{question}[Endpoint regularity beyond the Sobolev threshold]
Determine the complete Besov and logarithmic endpoint regularity of
the normalized Stieltjes primitives, including critical indices and
finite linear combinations.
\end{question}
The exact law $\|\Delta_hJ_n\|_2^2\sim
h\log^{2n+2}(1/h)/(n+1)^2$ already determines the first critical
obstruction. Higher finite differences and the explicit Fourier
polynomials should determine the remaining scales. A useful outcome
would specify convergence rates for Fourier and spline approximations
without assuming bounded endpoint derivatives.

\subsection{Uniform Euler limits}

\begin{question}[Effective reversal bounds]
Replace the asymptotic remainders in
Lemma~\ref{eutrans:lem:joint} by explicit constants on prescribed
$b$ intervals. Can the corrected Lambert expression be turned into
a certified bracket for $\nu_r$ at practical $\varepsilon$ values?
\end{question}
The proof already separates the error budgets. The task is to retain
numerical constants in the seed-moment, central-window, and tail bounds.
An exact positive-integral evaluation can then confirm signs at the
two proposed bracket endpoints. A small quadrature residual alone
would not certify the bracket.

\begin{question}[All-order and growing-parameter transitions]
Obtain an arbitrary-order expansion of the near-critical reversal
and identify the valid joint ranges when $r$, $b$, or $1/b$ grows
with $\lambda=\log(1/\varepsilon)$.
\end{question}
For fixed $r,b$, every further density moment and Gaussian logarithmic
moment is explicit in zeta and polygamma values. Growing $r$ moves the
Gaussian moment's mass, and $b\downarrow0$ changes the small-variable
seed bounds. These effects need a fresh localization. In particular,
the fixed-$r$ ratio $\nu_r/\nu_0\sim
[4^r/\binom{2r}{r}]^2$ cannot be used uniformly in growing $r$
without that analysis.

\begin{question}[Finite-index Euler extremizers]
Resolve the latest continuation's all-index axis and kernel-uniqueness
conjectures, beyond their already proved initial and eventual cases.
\end{question}
The question is the location of a global maximizer over both orders.
It is stronger than the eventual localization theorem, but different
from a pointwise comparison at fixed inner order. A practical route is
an effective eventual threshold plus exact algebraic or interval
analysis for the remaining finite indices. The new near-critical
reversal describes one potentially delicate boundary regime.

\subsection{Lerch phase arithmetic and critical sums}

\begin{question}[Phase arithmetic and quantitative sign statistics]
Is $\theta(2)/(2\pi)$ irrational? More generally, for algebraic $A>1$,
which arithmetic properties of $\theta(A)$ control the frequency and
length distribution of same-sign velocity runs?
\end{question}
The optimal block theorem needs no answer to this arithmetic question.
Sharper statistics do: equidistribution would give sign frequencies,
while quantitative bounds near the cosine zeros require both a
Diophantine estimate and the additive, rather than relative, coefficient
asymptotic. The reciprocal-angle cases also invite a classification of
the corresponding $A$ values.

\begin{question}[Uniform and smoothed critical finite parts]
Develop the critical-sum theorem uniformly as $A\downarrow1$ or
$A\to\infty$, and compare its sharp-cutoff constant with constants
obtained from exponential or smooth cutoffs.
\end{question}
The two singularities move with $A$, and the factor
$(1-e^{-2i\theta})^{-1}$ becomes large as the phase approaches an
endpoint. The fixed-$A$ theorem does not cover this coalescence.
Different cutoffs may retain the same local analytic jet but alter
subtracted terms; any equality of the resulting constants must be
proved from the corresponding generating operator.

\begin{question}[Global indexing between root regimes]
Prove the indexed bulk phase formula proposed in the latest Lerch
continuation, with its exact half-integer shift, and match the bulk
zeros to the small-root Bessel and large-root Gamma regimes.
\end{question}
The additive cosine expansion locates a local root near a phase value.
It does not on its own identify that phase with the globally ordered
$j$th zero. A uniform overlap argument with endpoint asymptotics is
needed to fix the index. This is also a route to discrepancy estimates
for the existing weak zero distribution.

\subsection{Exact polylogarithmic reductions and formalization}

\begin{question}[An exact certificate for the frozen mixed $S_6$ target]
After imposing the complete Cayley ideal with
$\Pi_{\uparrow}$, does the 30-term target lie in the span generated
by a precisely stated double-shuffle, distribution, and boundary
relation system? If so, produce an exact rational certificate.
\end{question}
At weight seven and depth at most four the ambient projected problem
has dimension $1964$. This makes a finite calculation more focused,
but no dimension count ensures that the target vanishes there.
If that specified system fails, the output should be an explicit
linear functional separating the target from that system, not a
claim of numerical nonreducibility. The current $S_8$ candidate needs
the same distinction between exact relation and tiny residual.

\begin{question}[Relations compatible with the oriented depth basis]
Construct effective multiplication and change-of-basis algorithms in
the filtered Cayley quotient, and determine which additional exact
polylogarithm relations admit similarly compatible depth projections.
\end{question}
The generator theorem gives existence and minimum formal depth. The
remaining algorithmic question is complexity at large weights, where
the raw Eulerian permutation sum is expensive. Lyndon recursion or
consecutive-cut dynamic programming may provide practical substitutes.
Any extension should specify its ideal before claiming completeness.

\begin{question}[Certified evaluation and proof-assistant targets]
Formalize the exact finite identities first, and then connect them to
the analytic theorems through a small collection of reusable lemmas.
\end{question}
Natural finite targets are the endpoint retraction, shuffle projector,
character counts, and polynomial coefficient extraction. Analytic
targets include the $L^2$ holomorphic Fourier formula, the removable
Laurent products, the epsilon-sensitive density estimate, and transfer
from a finite collection of slit singularities. The supplied scripts
are reproducible checks and prospective formalization inputs, rather
than completed proof-assistant developments.

\section{Reproducibility and verification record}
\label{sec:verification}

The package includes the modular article sources, an assembled single
TeX file, the rendered PDF, exact-algebra programs, high-precision
diagnostics, source provenance, and integration notes. Its README
gives complete commands and separates the quick exact checks from
the slower integral diagnostics. A SHA-256 manifest identifies the
delivered files; rebuilding a PDF may change its metadata and therefore
its binary hash.

\paragraph{Analytic identities.}
The translation program compares ordinary split Gamma quadrature
against the Hurwitz representation at several shifts, including a
small shift. It checks the two-order germ at real and complex orders,
the odd-zeta series with its analytic tail bound, the accelerated
$\gamma_1$ formula, and the second derivative. Rational-shift
polylogarithm order derivatives are evaluated independently through a
finite Hurwitz Fourier transform. Exact symbolic truncation constructs
the endpoint polynomials, rather than fitting their coefficients.
The supplement checks the Gram integral and the explicit primitive.

\paragraph{Euler transition.}
Nine samples evaluate the exact compensated kernel for $b=1,2$ and
derivative orders $r=0,1,2$. They compare root positions with the
leading, corrected, and Lambert approximations. A tighter-tolerance
repeat changed one logarithmic root by approximately
$4.1\times10^{-12}$. This is useful numerical evidence of stability;
the quadrature is not outward-rounded interval arithmetic. The
asymptotic theorem rests on the uniform proof, not these samples.

\paragraph{Lerch sums.}
The critical-sum program derives Puiseux coefficients symbolically,
compares two independent velocity recurrences at high precision, and
checks the $p=2,3,4$ subtractions on long finite sums. It also displays
the growing oscillation left by omitting the conjugate term.
The $O(N^{-1/2})$ theorem comes from the proved singular expansion
and continuation, not an extrapolation of the diagnostic errors.

\paragraph{Cayley algebra.}
All acceptance calculations in the projector suite use integers and
rational fractions. It tests $780$ words and $2150$ products, compares
two Eulerian formulas, and independently builds the entire ideal in
weights at most four. The independently obtained ranks at every depth
agree with the character formulas. The full frozen $S_6$ vector is
included with its rational normal form and exact nonzero coefficient.
These finite tests check the implementation; the all-weight
conclusions follow from the polynomial-algebra proofs.

\paragraph{Review status.}
The delicate Euler remainder, root inversion, and global maximum
identification received a separate analytic audit. The Cayley kernel
and exact minimum-depth argument received an independent algebraic
review. The final translation Gram and primitive identities were also
checked independently. These reviews are documented in the package.
The report was prepared with AI assistance and has not undergone
external peer review. Its claims are ordinary mathematical proofs
with stated dependencies and reproducible evidence, ready for the
project's integration and further mathematical review.
